\chapter{用户与对象理解}
\label{chap:rs-understanding}

某位用户昨天浏览了登山包、水壶和相机，今天要求“防水通勤背包，价格不超过500元”。系统已经有了一些记录，却还没有得到可以直接推荐的答案：昨天的相机浏览是否与今天有关，商品标题中的“户外”是否意味着适合通勤，未注明重量的背包是否满足轻便要求，都需要分别判断。本章把这些判断所需的信息组织起来，说明怎样从已知证据形成需求、对象及两者匹配的表示。

读者已经熟悉向量表示、监督学习以及循环网络、注意力和图网络。本章关心这些工具在推荐中接收什么证据、用什么监督学习、最终输出什么。相同历史可以被压成一个状态、保留为多个兴趣，也可以等候选到来后再选择相关部分；三种做法承担的计算接口不同。信息不足时，还需决定能否补充内容、迁移已有知识或取得新反馈。只有把这些条件说清楚，后续召回和排序才能知道自己实际拿到了什么依据。

% Outline: 2.1
\section{理解任务与信息来源}
\label{sec:rs-understanding-evidence}

理解环节提供三类相互关联的输出。用户侧给出需求与偏好依据，例如当前用途、价格上限和近期摄影行为；对象侧给出属性与关系，例如容量、重量、接口兼容性及供货状态；匹配侧把需求对应到对象，例如“价格合格、重量未知、通勤适用性有待核验”。匹配结果可以是结构化条件，也可以是服务某项预测的向量或分数。向量中的某一维没有被定义为重量时，就不能将其数值当作重量检查结果。

证据需要保留来源。用户亲自填写“不超过500元”是直接表达；日志记录了相机点击，是交互观测；模型由此推测“可能需要摄影用品”，则是推断。直接表达也可能含混或过时，但不能仅因历史模型给高价商品打了高分就被悄悄覆盖。交互观测还受到机会影响：用户点击了看到的商品，却未必有机会比较未展示的商品。因此，训练一个行为预测器与证明某项偏好成立之间仍有距离。

用户之间的关注、好友和信任关系可以提供额外依据，但它们与交互相似性不是同一关系。两人互相关注，不意味着购买过相同物品；两人买过相同相机，也不证明存在社交联系。对象之间的品牌、部件和主题关系又属于另一类图。第~\ref{chap:rs-retrieval}~章会利用这些关系形成目录候选；本章先要求记录清楚边的含义和取得时间，避免不同关系被统一写成不明来源的“相似”。

一次请求也可能属于一个群组。共同选电影时，应知道哪些成员参加、哪个反馈属于共同选择，不能把某位成员事后的个人观看混入群组标签。个人与群组模型都需要在请求时点$t$截断信息。以下用$H_{u,t}=(e_1,\ldots,e_n)$表示已知历史，事件$e_k$至少包含物品$i_k$、动作$b_k$及时间$\tau_k$。浏览、收藏、加购与购买可以共同作为输入，但它们的标签和预测窗口须分别定义。

贯穿本章的商品例子只用于检查计算。我们会明确给出假设向量、权重和时间，不把算出的兴趣表示当作用户真实心理，也不把构造数据当作论文实验。固定相同事件再改变编码方式，才能看清方法差别来自顺序、候选选择还是额外信息。

% Outline: 2.2
\section{需求与对象建模}
\label{sec:rs-understanding-models}

需求、对象和匹配是本节的三个观察角度。用户模型回答哪些历史或表达应影响当前判断，对象模型回答有哪些可用事实及关系，匹配模型回答这些信息怎样共同支持选择。实际网络可以联合训练三者，不必划成三个独立服务；但讲解时仍要分清中间表示与最终预测，才能判断某个模块是否可以移到召回、排序或其他任务中。

% Outline: 2.2.1
\subsection{用户需求理解}
\label{sec:rs-user-understanding}

当前表达通常比一份长期画像更贴近这次任务。“只看500元以内”是明确限制，“轻便一些”需要结合场景解释，“不要这款”还需知道拒绝的是价格、款式还是已经购买。时间、地点、设备及刚刚进行的操作也会改变需求。历史可以填补未表达的部分，遇到冲突时则应保留表达的优先级或发起澄清。查询怎样对应搜索相关性，见第~\ref{chap:rs-search}~章。

行为建模面对的几组选择不能混为互斥类别。一个模型可以同时使用长历史、多兴趣表示和候选相关聚合。单一表示与多兴趣表示比较的是保留多少个兴趣方向，候选无关与候选相关比较的是候选在哪一步进入计算，短会话与长历史比较的是可用证据范围。例如，登山与摄影可以分别保留两个向量，再由多个向量取得候选；也可以保留完整历史，在评分相机时选择摄影记录。前一种目录接口在第~\ref{chap:rs-retrieval}~章的多兴趣模型中展开。

训练目标同样不能由输入形式推断。使用购买和收藏历史，不表示两个事件具有相同价值；预测下一次浏览，也不同于预测给定曝光是否点击。表~\ref{tab:rs-behavior-models}先比较常用模型的完整接口。辅助行为怎样支持多个预测任务，见第~\ref{chap:rs-ranking}~章；这些预测目标如何服务业务取舍，见第~\ref{chap:rs-constraints}~章。

\begin{table*}[htbp]
  \centering
  \small
  \renewcommand{\arraystretch}{1.15}
  \begin{tabularx}{\linewidth}{@{}lXXXX@{}}
    \toprule
    模型 & 输入及组织 & 候选参与位置 & 中间输出 & 监督与应用输出 \\
    \midrule
    GRU4Rec & 匿名会话，逐事件更新 & 状态编码后 & 会话循环状态 & 下一物品比较；目录得分 \\
    SASRec & 物品前缀，因果注意 & 历史编码后 & 末端历史向量 & 后继正负例；目录分数 \\
    BERT4Rec & 已知历史，掩码重建 & 历史编码后 & MASK位置向量 & 遮蔽物品；目录概率 \\
    SR-GNN & 会话内转移图 & 会话编码后 & 局部与整体会话向量 & 下一点击；目录得分 \\
    DIN & 行为集合 & 历史激活 & 候选相关兴趣 & 曝光点击；CTR \\
    DIEN & 顺序历史及辅助目标 & 第二层兴趣演化 & 候选相关循环状态 & 下一行为及点击；CTR \\
    DSIN & 划分后的历史会话 & 会话级聚合 & 两组会话兴趣 & 曝光点击；CTR \\
    MIMN & 增量行为、固定记忆 & 预测网络 & 记忆及归纳状态 & 点击与记忆正则；CTR \\
    SIM & 完整历史与检索子集 & 检索及精细聚合 & 相关历史兴趣 & 模块及最终点击；CTR \\
    TWIN & 长历史、共享相关性 & 两阶段选择 & 多头目标兴趣 & 曝光点击；CTR \\
    TransAct & 即时动作、批量表示 & 逐行为早融合 & 压缩后的近期表示 & 多行为标签；响应分数 \\
    HSTU & 物品与动作异质序列 & 随召回／排序任务变化 & 共享序列状态 & 召回分布或动作概率 \\
    \bottomrule
  \end{tabularx}
  \caption{行为模型的证据、表示和完整使用接口。目录得分可供后续候选取得，CTR则需要一个给定候选；同样使用历史序列，不代表两种接口能够原样互换。}
  \label{tab:rs-behavior-models}
\end{table*}

% Outline: 2.2.1.1
\subsubsection{GRU4Rec：循环状态的会话建模}
\label{sec:rs-gru4rec}

匿名访问没有可用的长期用户向量，但当前会话中的连续点击仍然提供了线索。\term{GRU4Rec}用门控循环单元（GRU）随事件更新会话状态，再为下一物品计算得分。原模型以物品标识的独热向量作为输入；论文也考察了历史独热向量的衰减聚合和附加嵌入层，不能将后来常见的“先查嵌入再输入GRU”说成其唯一原始设计\scite{rs-r10}

设$\bm{x}_k\in\{0,1\}^{|I|}$对应第$k$次点击，隐藏状态为$\bm{h}_k\in\mathbb{R}^d$。为了看清新事件怎样改变依据，可把GRU更新写成
\begin{equation}
  \bm{h}_k=(\bm{1}-\bm{z}_k)\odot\bm{h}_{k-1}
       +\bm{z}_k\odot\widetilde{\bm{h}}_k,\qquad
  s_{k,i}=f_{\mathrm{out},i}(\bm{h}_k),
  \label{eq:rs-gru4rec-state}
\end{equation}
其中更新门$\bm{z}_k$和提议状态$\widetilde{\bm{h}}_k$由当前独热输入、此前状态及重置门共同计算，$f_{\mathrm{out}}$是输出层。这里采用“门越大，新状态写入越多”的记号；另一种GRU约定可把门取补，含义不变。

训练要让实际下一物品$i^+$超过采样物品$j$。令$J_k$为本步负例集合，$\sigma$为sigmoid，原文讨论的两项损失为
\begin{align}
  \ell_{\mathrm{BPR}}
    &=-\frac{1}{|J_k|}\sum_{j\in J_k}\log\sigma(s_{k,i^+}-s_{k,j}),\\
  \ell_{\mathrm{TOP1}}
    &=\frac{1}{|J_k|}\sum_{j\in J_k}
       \bigl[\sigma(s_{k,j}-s_{k,i^+})+\sigma(s_{k,j}^2)\bigr].
  \label{eq:rs-gru4rec-top1}
\end{align}
TOP1的第一项平滑近似负例排在正例前的程度，第二项把负例得分向零约束，避免只靠整体抬分降低相对比较项。未观测物品在这里是训练对照，不是已证实不喜欢的物品。

会话并行批次在每个槽位放一段会话，依次处理各会话的下一事件；某个会话结束后，换入新会话并重置该槽位的状态。原方法利用同批其他样本的目标物品作为负例，节省输出计算，但也使采样频率与物品流行度相关。推断时逐点击更新状态，计算目录或候选上的得分。

例如，假设点击登山包后的状态为$(0.6,0.1)$，水壶事件产生提议状态$(0.8,0.2)$和门$(0.5,0.5)$，则新状态为$(0.7,0.15)$。若教学输出层对水壶配件与相机分别采用向量$(1,0)$和$(0,1)$，得分就是0.7和0.15。状态改变了下一次比较，却没有建立跨会话记忆：新会话重置后，这段个体信息不会自行保留。

% Outline: 2.2.1.2
\subsubsection{SASRec：自注意力的历史序列建模}
\label{sec:rs-sasrec}

循环状态把全部前缀压在一个递推变量中；若登山与摄影事件交错，预测时也可以直接读取此前位置。\term{SASRec}将物品嵌入与位置嵌入相加，以因果自注意力编码前缀，再用末端表示与物品向量比较。候选物品参与最后的评分，不进入历史编码本身\scite{rs-r11}

设截断后的历史长$n$，输入矩阵为$\bm{X}\in\mathbb{R}^{n\times d}$，其第$k$行为$\bm{v}_{i_k}+\bm{p}_k$。一层注意力的核心计算是
\begin{equation}
  \bm{A}=\operatorname{softmax}_{\mathrm{row}}
  \left(\frac{\bm{Q}\bm{K}^{\top}}{\sqrt d}+\bm{M}\right),
  \qquad \bm{Z}=\bm{A}\bm{V},
  \label{eq:rs-sasrec-attention}
\end{equation}
其中$\bm{Q},\bm{K},\bm{V}$由输入线性投影得到。掩码$\bm{M}$把未来位置及填充位置排除；填充不应因为恰好有零向量而获得一份注意权重。位置前馈变换增加非线性，残差保留原表示，层归一化和dropout支持训练。通用运算见\BookChapterRef{chap:transformer}{模型卷的“Transformer”章}。

训练时每个合法位置$k$用$i_{k+1}$作正例，并从目录采样$j_k$作对照，最小化
\begin{equation}
  \ell_k=-\log\sigma(\langle\bm{h}_k,\bm{v}_{i_{k+1}}\rangle)
  -\log\bigl[1-\sigma(\langle\bm{h}_k,\bm{v}_{j_k}\rangle)\bigr].
  \label{eq:rs-sasrec-loss}
\end{equation}
推断取$\bm{h}_n$，为目录物品计算内积。训练位置$k$不能看到$i_{k+1}$，即使这个物品已存在于整条训练序列中；因果掩码正是为了让共享计算仍符合各位置的预测条件。

假设历史为登山包、水壶、相机，最后一行注意权重为$(0.2,0.1,0.7)$，值向量依次为$(1,0),(0.8,0.2),(0,1)$，则该层聚合结果为$(0.28,0.72)$。这一计算说明末端可以保留多个先前位置的影响，不能据此宣称任意相机候选都会再次改变该组权重。长度为$n$的稠密注意力仍需计算约$n^2$个位置关系；截断会遗漏窗口外事件，扩窗则增加成本。完整目录排名和少量采样候选排名也不是同一难度，评价时须分别报告。

% Outline: 2.2.1.3
\subsubsection{BERT4Rec：掩码训练的双向历史建模}
\label{sec:rs-bert4rec}

下一物品目标只从左侧前缀取得监督。已经完整观察到的一段历史还可以提供另一种练习：遮住其中某个物品，让模型根据其余已知事件恢复它。\term{BERT4Rec}采用这种掩码物品重建，以双向Transformer编码物品与位置嵌入。它借鉴语言建模任务的形式，但编码的词表是物品目录，并非直接把NLP中预训练好的BERT当作推荐器\scite{rs-r26}

令$S$为被选作重建目标的位置集合，$\widetilde H$为经过遮蔽或替换的输入，$\bm{h}_k$为其第$k$个输出。通过变换后的$\bm{h}_k$和物品输出嵌入得到目录概率，训练损失为
\begin{equation}
  \ell_{\mathrm{mask}}
    =-\frac{1}{|S|}\sum_{k\in S}
        \log p_\theta(i_k\mid\widetilde H).
  \label{eq:rs-bert4rec-mask}
\end{equation}
由于目标物品已经被隐藏，内部位置可以读取左右两侧的其他已知物品。推断下一物品时，在请求前历史末端追加MASK，以该位置输出分布；训练还加入末项遮蔽情形，使模型接触这一使用接口。

对“登山包、MASK、相机”的重建，答案可能是曾被遮蔽的水壶；对“登山包、水壶、相机、MASK”的推断，目标则是尚未发生的下一行为。假设后一种输入给镜头、水壶、背包配件的概率为0.5、0.2、0.3，则本次排序为镜头、配件、水壶。前一个训练问题的答案不能直接搬到后一个问题上。

双向的范围是请求前已知历史。训练切分、掩码位置及末项构造一旦把请求后的行为混入，就改变了可用证据。完整目录softmax还可能昂贵；减少训练候选需要重新说明目标近似和评价口径。模型擅长何种历史条件，应由同一信息边界下的比较确定。

% Outline: 2.2.1.4
\subsubsection{SR-GNN：物品转移图的会话建模}
\label{sec:rs-srgnn}

重复访问会让会话既包含顺序，也包含回访和分支。\term{SR-GNN}将一次会话的不同物品建成节点，以相邻点击构造有向转移边；通过门控图传播整合入边与出边信息，再组合末次物品和整个会话的表示，预测下一次点击\scite{rs-r13}

设节点$i$的表示为$\bm{v}_i^{(\ell)}$，$\bm{A}^{\mathrm{in}}$和$\bm{A}^{\mathrm{out}}$按相应连接关系归一化。一次传播先分别汇入前驱和后继表示，将两侧线性变换的结果拼接成$\bm{a}_i^{(\ell)}$，再以GRU式门更新节点。重复转移须按固定规则计数或合并，并在相应出入邻域内归一化；同一规则在训练和服务中保持一致。

末次点击$i_n$给出局部表示$\bm{s}_{\mathrm{local}}=\bm{v}_{i_n}^{(L)}$，会话整体表示为
\begin{equation}
  \bm{s}_{\mathrm{global}}=\sum_{k=1}^{n}\alpha_k\bm{v}_{i_k}^{(L)},
  \qquad
  \bm{s}=\bm{W}
  [\bm{s}_{\mathrm{local}};\bm{s}_{\mathrm{global}}],
  \label{eq:rs-srgnn-session}
\end{equation}
其中$\alpha_k=\bm{q}^{\top}\sigma(\bm{W}_1\bm{v}_{i_n}^{(L)}
+\bm{W}_2\bm{v}_{i_k}^{(L)}+\bm{c})$由最后事件和各事件共同计算。原形式的$\alpha_k$不是必须归一化到和为一的概率。用$\langle\bm{s},\bm{v}_i\rangle$形成物品logit，对实际下一点击的softmax交叉熵进行训练。

会话“登山包、水壶、登山包、相机”产生三条转移：包到水壶、水壶到包、包到相机。包的两个不同后继在等权情况下各占出边归一化权重的$1/2$；末次节点是相机，整体聚合中登山包又出现两次。若再发生一次包到水壶，采用转移频次的图会改变这两个后继的比例。图因此能表达回访结构，但只保留节点、边和用于聚合的事件信息，不能从邻接矩阵唯一还原完整事件顺序。它也没有包含其他用户的协同关系。

% Outline: 2.2.1.5
\subsubsection{DIN：候选相关的历史兴趣聚合}
\label{sec:rs-din}

前面的模型先形成一份会话或历史表示，再与各物品比较。预估给定广告的点击时，历史中哪些事件有用可以直接由广告决定。\term{DIN}的局部激活单元逐条比较候选与行为，形成候选相关的兴趣向量，再与其他字段共同预测CTR，即给定曝光条件下点击发生的概率\scite{rs-r12}

以候选向量$\bm{v}_i$和行为向量$\bm{v}_{i_k}$为输入，原局部激活单元利用两者及其交互信息。可将其写为
\begin{equation}
  a_k(i)=g_\theta([\bm{v}_{i_k};\bm{v}_i;
            \bm{v}_{i_k}-\bm{v}_i;
            \bm{v}_{i_k}\odot\bm{v}_i]),\qquad
  \bm{h}(i)=\sum_{k=1}^{n}a_k(i)\bm{v}_{i_k}.
  \label{eq:rs-din-activation}
\end{equation}
权重无需softmax归一化，因此相同相关事件的数量也可能改变表示幅度。将$\bm{h}(i)$、用户字段、候选及情境送入预测网络，得到$\widehat p_i=\sigma(f_\theta(\cdot))$，用曝光点击二元交叉熵联合更新嵌入、激活单元和预测层。

假设登山包、水壶、相机的向量为$(1,0),(0.8,0.2),(0,1)$。对背包候选，教学权重为$(0.8,0.6,0.1)$，聚合得到$(1.28,0.22)$；对相机候选，权重为$(0.1,0.1,0.9)$，则得到$(0.18,0.92)$。如果仅为演示后续预测而取$f(\bm{h})=h_1+h_2-1$，两者CTR为$\sigma(0.50)$和$\sigma(0.10)$，约0.622和0.525。向量差异与最终概率排序由不同计算步骤决定，不能只看某个兴趣维度作结论。

原模型还使用Dice激活函数，按批量统计量形成数据依赖的门，在两条斜率分支间平滑选择；小批量感知正则化按稀疏特征出现情况分摊嵌入惩罚，减少每步遍历整张嵌入表的需求。这些设计与候选相关聚合共同训练，但各自解决非线性适配和稀疏正则计算问题。

DIN的求和本身不恢复完整顺序，CTR也依赖曝光样本和标签定义。候选变了要重新计算激活，因而这份兴趣向量不能原样缓存成一份候选无关用户塔。注意权重解释的是网络采用了哪些输入，不足以证明真实动机，更不能直接当作所有后续目标的价值。

% Outline: 2.2.1.6
\subsubsection{DIEN：候选相关的兴趣演化}
\label{sec:rs-dien}

DIN可以选择相关事件，却未明确区分一串事件中的中间兴趣状态。\term{DIEN}先用GRU提取每一步状态，再让这些状态依据候选参与第二层演化。两层的分工是：第一层学习能解释随后行为的历史表示，第二层决定哪些变化与当前候选相关\scite{rs-r24}

仅靠会话末端的点击标签，早期状态很难得到直接约束。DIEN为第$k$步状态$\bm{h}_k$增加下一行为辅助任务，区分真实$i_{k+1}$与采样$\widetilde i_{k+1}$：
\begin{equation}
  \ell_{\mathrm{aux},k}
  =-\log D(\bm{h}_k,\bm{v}_{i_{k+1}})
   -\log[1-D(\bm{h}_k,\bm{v}_{\widetilde i_{k+1}})].
  \label{eq:rs-dien-aux}
\end{equation}
$D$是辅助判别网络。下一行为只用作监督，不能进入生成$\bm{h}_k$的前缀输入。实际CTR损失与辅助损失按系数$\lambda$相加，使中间状态和最终任务一起学习。

第二层先计算候选与$\bm{h}_k$的注意权重$a_k(i)$，再通过注意更新门GRU（AUGRU）调节写入幅度。若普通更新门为$\bm{z}_k$，则
\begin{equation}
  \widetilde{\bm{z}}_k=a_k(i)\bm{z}_k,\qquad
  \bm{g}_k=(\bm{1}-\widetilde{\bm{z}}_k)\odot\bm{g}_{k-1}
               +\widetilde{\bm{z}}_k\odot\widetilde{\bm{g}}_k.
  \label{eq:rs-dien-augru}
\end{equation}
AIGRU把注意权重乘到输入，弱输入仍可能触发状态变化；AGRU用标量注意力直接替换更新门；AUGRU保留各维门值，再整体缩放。假设$\bm{z}_k=(0.8,0.2)$、$a_k=0.25$，AUGRU得到$(0.2,0.05)$，而AGRU在两维都使用0.25，二者改变的计算不同。

在登山与摄影交错历史中，相机候选可以降低与登山状态对应的写入，让摄影相关轨迹更多保留到末端。最终$\bm{g}_n$与用户、候选和情境送入CTR网络。模型输出的是受点击及辅助任务约束的预测状态；称它为“兴趣演化”，并不意味着已经识别真实心理变化，更没有比较推荐行动怎样改变长期结果。

% Outline: 2.2.1.7
\subsubsection{DSIN：会话内兴趣与跨会话关系}
\label{sec:rs-dsin}

行为按时间紧邻，并不一定来自同一个连续任务。上午集中选登山用品、晚上集中看摄影器材时，直接把全部事件看作一串等价位置可能掩盖会话结构。\term{DSIN}先按时间规则划分会话，在会话内提取兴趣，再建模会话之间的联系\scite{rs-r25}

设历史分为$m$个会话$H^{(1)},\ldots,H^{(m)}$。每个会话的行为嵌入加入区分会话、位置及特征维度的偏置编码，通过多头自注意力和池化形成$\bm{s}_r$。这些偏置告诉网络同一嵌入位于哪里；会话边界本身则由预处理规则给定，不是注意力自动发现的真实意图切换。双向LSTM继续读取请求前的全部会话向量，形成$\bm{g}_r$。

候选对两组向量分别进行激活：
\begin{equation}
  \bm{h}_{\mathrm{session}}(i)=\sum_r\alpha_r(i)\bm{s}_r,\qquad
  \bm{h}_{\mathrm{relation}}(i)=\sum_r\beta_r(i)\bm{g}_r.
  \label{eq:rs-dsin-pooling}
\end{equation}
两组归一化权重由候选和相应会话表示的兼容程度决定。预测网络结合这两个输出及其他用户、物品字段，以曝光点击二元交叉熵训练全部可学习模块。

假设两个会话的向量分别为$(1,0)$与$(0,1)$，相机候选给它们的激活权重为0.2和0.8，则第一组聚合为$(0.2,0.8)$；第二组在双向LSTM之后计算，已经包含跨会话关系，不能直接复用这个数值。服务时同一候选同时读取两组聚合，而不是仅取最后一个会话。

双向计算仅覆盖预测时已知的历史会话。若会话划分把两个用途混在一起，或者因时间阈值把同一任务截断，模型建立的关系也会改变。因而应在相同事件上比较划分规则，并检验短会话、混合任务和边界附近的表现，而不能将“会话内部兴趣一致”当作全部数据都成立的事实。

% Outline: 2.2.1.8
\subsubsection{MIMN：增量记忆的长历史压缩}
\label{sec:rs-mimn}

每次请求都重新扫描完整历史，成本会随事件数持续增加。\term{MIMN}把一部分工作移到行为到来时：维护固定数量的记忆槽，新事件读取和修改记忆，预测请求再使用已经维护的状态。多通道用户兴趣记忆网络由基础读写记忆和记忆归纳单元共同组成\scite{rs-r27}

设记忆矩阵$\bm{M}^{(k)}\in\mathbb{R}^{m\times d}$包含$m$个槽，控制器由新行为生成读取键、写入键、擦除向量$\bm{e}_k$和添加向量$\bm{a}_k$。内容寻址将键与各槽的余弦相似度归一化为读取或写入权重。用$\bm{w}^{\mathrm{r}}_k,\bm{w}^{\mathrm{w}}_k\in\mathbb{R}^m$表示这些权重，则基本操作为
\begin{align}
  \bm{r}_k&=\sum_{j=1}^{m}w^{\mathrm{r}}_{k,j}\bm{M}^{(k-1)}_{j,:},\\
  \bm{M}^{(k)}_{j,:}
    &=\bm{M}^{(k-1)}_{j,:}\odot
        (\bm{1}-w^{\mathrm{w}}_{k,j}\bm{e}_k)
       +w^{\mathrm{w}}_{k,j}\bm{a}_k.
  \label{eq:rs-mimn-write}
\end{align}
写入因此既可减弱旧内容，也可加入新信息。模型还根据累计使用量重新平衡写入权重，并惩罚各槽累计写入量的方差，以减少少数槽长期承受几乎全部更新的情况。

单纯槽值还不足以表示更新轨迹。记忆归纳单元为槽位维护多通道GRU状态$\bm{S}$，选择读取权重较高的通道，结合相应记忆与当前行为更新；各通道共享GRU参数。读出及归纳状态与候选、用户和情境一起进入CTR网络，训练由点击损失和记忆使用正则共同约束。

假设两个槽原为$(1,0)$与$(0,1)$，新摄影收藏的写入权重为$(0.2,0.8)$，擦除向量为$(0.5,0.5)$，添加向量为$(0,1)$。按式~\eqref{eq:rs-mimn-write}，新槽为$(0.9,0.2)$与$(0,1.4)$。同一事件改变了两个槽，幅度不同；槽的含义来自学习，不能预先断言每个槽始终对应一个固定兴趣类别。

原方案中的用户兴趣中心把这些更新放在行为触发路径，服务请求读取状态，减少重复历史计算。固定$m$意味着在线读写规模不直接随已发生事件总数增长，但状态存储、行为更新和版本同步仍有代价，历史细节也可能丢失。新行为是否已经写入，属于第~\ref{sec:rs-evidence-refresh}~节的时效问题；能够持续写入多年历史，不等于无损保存多年事件。

% Outline: 2.2.1.9
\subsubsection{SIM：长历史的检索与精细匹配}
\label{sec:rs-sim}

固定容量记忆节省了重复读取，却难以恢复已经压掉的事件细节。\term{SIM}采用另一条路径：保留原始长历史，先按候选寻找有关事件，再精细处理选出的子序列。通用检索单元（general search unit, GSU）负责缩小历史范围，精细匹配单元（exact search unit, ESU）负责聚合；这里检索的是用户过去的行为，不是待推荐的目录物品\scite{rs-r20}

硬检索用类别等规则确定关联，例如相机候选访问用户摄影类历史的倒排记录。软检索则学习物品表示，用内积等度量取得最相关的$K$条行为，原方案讨论以近似最大内积检索支持该步骤。令$S_i\subseteq\{1,\ldots,n\}$为检索出的历史位置，精细匹配把行为与距请求的时间间隔编码为$\bm{e}_k$，计算
\begin{equation}
  \bm{h}(i)=\mathop{\Vert}_{r=1}^{h}
   \sum_{k\in S_i}\alpha_{r,k}(i)\bm{W}^{\mathrm{v}}_r\bm{e}_k,
  \label{eq:rs-sim-interest}
\end{equation}
其中$h$为注意头数，$\Vert$表示向量拼接，各头权重由投影后的候选和选中行为计算。该表示与短期兴趣及其他特征共同送入CTR网络。

软GSU有基于长历史聚合的辅助点击预测，ESU有最终点击预测，两项交叉熵加权训练并共享相关嵌入；硬GSU没有可学习的检索参数，因此不需要同样的GSU损失。可学习表示接受点击监督，不意味着离散Top-$K$操作已经成为普通可微层。训练及服务都须说明选集怎样产生、表示和索引何时同步。

假设六条历史依次为登山包、水壶、镜头、帐篷、相机、存储卡。摄影类别硬检索得到第3、5、6条；预算仅允许两条时，还要定义二次取舍。若取相机、存储卡，ESU就无法再利用镜头记录，即使完整历史中镜头最能解释当前需求。软检索可以改变候选子集，但同样存在漏检。评价必须共同固定原始历史范围、选取条数和计算预算，不能只以“用了多年历史”断言理解更好。

% Outline: 2.2.1.10
\subsubsection{TWIN：两阶段历史选择的相关性一致性}
\label{sec:rs-twin}

两阶段方法还有一个接口问题：粗阶段按类别或内积选出的事件，可能与精细阶段实际重视的事件不同。\term{TWIN}让两阶段采用相同的候选—行为相关性计算，通过拆分特征降低长历史扫描成本。这样做的目标是让有限的精细计算留给同一度量认为相关的事件\scite{rs-r28}

物品自身特征的投影可以跨用户复用，点击时间、观看情况等用户—物品交叉特征则随事件变化。TWIN将前者投影并缓存，把后者各自压成低维偏置。省略注意头编号后，一条历史的相关性可概括为
\begin{equation}
  a_k(i)=
  \frac{\langle\bm{W}^{\mathrm{q}}\bm{v}_i,
                 \bm{W}^{\mathrm{h}}\bm{f}_{i_k}\rangle}{\sqrt d}
       +\bm{\beta}^{\top}\bm{c}_{u,k},
  \label{eq:rs-twin-score}
\end{equation}
其中$\bm{f}_{i_k}$为物品自身特征，$\bm{c}_{u,k}$为压缩后的事件交叉特征，$\bm{\beta}$为可学习权重。粗阶段据此选择历史，精细阶段在选集内对相同相关性作softmax，再聚合完整值投影；多个头的结果与近期行为等分支共同预测CTR。

假设相机候选对镜头、相机包和水壶的精细相关性为3、2、0，另一套粗度量却给出1、0、2。只保留两条时，后者漏掉相机包。统一度量保留镜头和相机包，精细权重为$e^3/(e^3+e^2)\approx0.731$和0.269。例子说明一致性解决了什么，也说明它没有证明第三条完全无用：Top-$K$仍然截断了尾部证据。

点击监督训练共享相关性与预测网络，服务端复用其参数和物品投影。参数共享、参数同步及投影缓存刷新是不同要求：两阶段公式相同，缓存过时仍可让实际得分不同。相关性本身不准时，一致地选错也无济于事；因此应同时检查选择重合程度、最终预测和缓存时效，而不是只报告一致率。

% Outline: 2.2.1.11
\subsubsection{TransAct：即时行为与长期表示的组合}
\label{sec:rs-transact}

用户刚刚收藏摄影内容，批量更新的长期表示可能要稍后才会反映这一变化。\term{TransAct}把即时动作作为额外输入，与批量长期表示共同参与排序。它并不要求先重训整个长期模型，而是直接让当前请求读取近期可用证据\scite{rs-r14}

每条近期事件包含物品内容向量、动作类型和时间戳。模型将动作嵌入与内容向量结合，再把候选向量拼接到每条行为上，形成输入$\bm{x}_k(i)=[\bm{v}_{i_k};\bm{e}_{b_k};\bm{v}_i]$。这种早融合不同于只把候选追加成序列末尾的一个位置：前者从每个输入位置就带有候选信息，后者依赖后续注意力交换。原工作比较后选择了逐行为拼接。

Transformer编码按时间倒序排列的近期输入，取最近$K$个输出，并拼接全序列的逐维最大池化：
\begin{equation}
  \bm{h}_{\mathrm{recent}}(i)=
    [\bm{o}_1;\ldots;\bm{o}_K;
      \operatorname{maxpool}(\bm{o}_1,\ldots,\bm{o}_n)].
  \label{eq:rs-transact-compress}
\end{equation}
所得紧凑表示与长期表示及其他字段一起进入排序预测网络，用不同动作标签的加权预测损失训练。原配置未采用位置编码；时间戳用于排序和时间窗口掩码。因此，不能仅凭“使用Transformer”就宣称网络完整保留了全部位置关系。

训练随机抽取一个时间窗口，屏蔽请求前该窗口内的近期事件，促使模型也利用较远证据；推断不使用这个随机时间窗口掩码。它不同于防止未来泄漏的因果掩码，二者目的和部署方式都不同。假设近期两个输出为$(0.2,0.8)$与$(0.7,0.3)$、$K=1$，压缩结果为$(0.2,0.8,0.7,0.8)$，近期重点和逐维最大值都进入下游。

刚到达的收藏能改变这个分支，即使长期向量仍是昨日版本。但事件是否已被服务端读到、长期向量何时刷新、预测参数多久更新，依然是三条不同时间线。近期窗口也会遗漏更早事件；即时性改善需要在同一请求和标签窗口下验证，不能直接推导长期收益。

% Outline: 2.2.1.12
\subsubsection{HSTU：异质行为的序列转导}
\label{sec:rs-hstu}

推荐输入既有物品，也有点击、收藏和跳过等动作；不同任务又反复读取同一历史。\term{HSTU}把这些信息组织为异质序列，研究怎样共享多个合法位置的序列计算。序列转导在这里表示从已知事件序列预测后续物品或候选动作，输出对象由任务定义，而不是自由生成一个商品名称\scite{rs-r07}

令输入矩阵为$\bm{X}\in\mathbb{R}^{n\times d}$。一个HSTU层用逐位置投影同时生成查询、键、值及门控张量，并以位置和时间偏置调整聚合。省略多头重排后可写成
\begin{align}
  (\bm{U},\bm{V},\bm{Q},\bm{K})
    &=\operatorname{Split}\bigl(\operatorname{SiLU}(f_1(\bm{X}))\bigr),\\
  \bm{A}&=\operatorname{SiLU}(\bm{Q}\bm{K}^{\top}+\bm{B}),\\
  \bm{Y}&=f_2\bigl(\operatorname{Norm}(\bm{A}\bm{V})\odot\bm{U}\bigr).
  \label{eq:rs-hstu}
\end{align}
可见性掩码排除不合法的事件关系，层间使用残差连接。$\bm{B}$编码相对位置与时间，$f_1,f_2$为逐位置映射。这里逐对应用SiLU，没有将整行权重softmax归一化。

这一差别影响证据数量的表达：若两个完全相同的值都获得相同softmax权重，它们的加权平均与只有一个值时相同；未作行归一化的聚合则可随重复证据增加幅度。HSTU随后仍有归一化、门控和投影，因此实际输出怎样保留强度要由整层计算判断，不能把未归一化直接解释为次数计数器。

召回训练在合法前缀位置预测下一物品，用户状态与物品输出表示构成目录分布，实际计算可采用采样目标。排序训练则把待判断物品放在相应输入位置，用当时尚未知晓的动作作标签，接入响应预测头。历史中的既有动作可作为输入，当前候选的待预测动作不能提前写入。一次浏览后收藏与浏览后跳过可以共享物品信息，却必须形成不同的动作监督。

例如，历史为“背包浏览、水壶收藏、相机跳过”，召回头回答下一件可能交互的物品，排序头回答给定新背包曝光后的某种动作概率。两者可以复用序列编码结构，但候选是否输入、损失位置和输出头都不同。推断时能复用多少历史计算，还取决于候选之间和候选对历史的可见性。稀疏训练、随机长度和M-FALCON等服务设计属于系统实现；论文在特定规模上的扩展结果不能被省去条件地推广到所有数据和预算。

% Outline: 2.2.1.13
\subsubsection{ETA：候选相关的高效长历史选择}
\label{sec:rs-eta}

长历史检索可以使用便宜的二进制比较，但检索向量仍应反映点击任务。\term{ETA}用CTR网络中的物品嵌入生成SimHash指纹，通过候选与行为指纹的汉明距离选择相关历史，再对选中事件作目标注意力，与短期分支共同预测点击\scite{rs-r76}

给定随机投影方向$\bm{r}_q$，第$q$位指纹为
\begin{equation}
  b_q(\bm{v})=\mathbb{I}\{\langle\bm{r}_q,\bm{v}\rangle\geq0\},
  \qquad
  d_{\mathrm{H}}(\bm{b}_i,\bm{b}_{i_k})
     =\sum_{q=1}^{m}\mathbb{I}\{b_{i,q}\ne b_{i_k,q}\}.
  \label{eq:rs-eta-hash}
\end{equation}
在各向同性随机超平面和非零向量下，单个位的不同概率与向量夹角成正比，因此汉明距离可近似反映方向差异。选出的$K$条行为进入精细注意力；点击损失更新共享嵌入及预测网络，更新后的嵌入又会影响下一次指纹构造。

假设背包候选指纹为1101，三条历史为1101、1001、0010，则距离为0、1、4。保留两条时，前两条进入精细计算，第三条被排除。若第三条的向量恰好因有限随机投影而被误分，后续注意力无法补回。共享表示缓解了检索与评分各用一套训练依据的问题，但离散符号和Top-$K$不因此成为处处可微运算。指纹更新、长历史扫描与精细计算都需要计入成本。

% Outline: 2.2.1.14
\subsubsection{SDIM：多哈希的长历史聚合}
\label{sec:rs-sdim}

ETA找到历史后仍逐条计算注意力。候选数量很大时，可以进一步把候选相关聚合转成桶访问。\term{SDIM}进行多轮随机投影，每轮用若干位组成签名，将与候选同签名的行为向量相加，归一化后跨轮平均，形成长历史兴趣\scite{rs-r77}

设第$r$轮使用$b$位签名$h_r$，相应桶向量为
\begin{equation}
  \bm{B}_{r,c}=\sum_{k:h_r(\bm{v}_{i_k})=c}\bm{v}_{i_k},\qquad
  \bm{h}(i)=\frac{1}{R}\sum_{r=1}^{R}
       \operatorname{norm}_2(\bm{B}_{r,h_r(\bm{v}_i)}).
  \label{eq:rs-sdim-buckets}
\end{equation}
这里$\operatorname{norm}_2(\bm{z})=\bm{z}/\lVert\bm{z}\rVert_2$，空桶按零表示处理。桶内聚合可在行为侧维护，候选只需计算签名并取得相应桶表示。长历史兴趣与短期兴趣和其他特征一起接入CTR网络，用点击交叉熵训练。

若两个非零向量夹角为$\theta$，独立随机超平面下，$b$位全部碰撞的概率为$(1-\theta/\pi)^b$。夹角为$60^\circ$时，两位签名的碰撞概率为$4/9$，四位降为$16/81$。签名更宽减少不同方向混入，也会漏掉更多有一定相似性的事件；增加轮数可改变这种取舍。假设某轮同桶向量为$(1,0)$和$(0.8,0.2)$，其和$(1.8,0.2)$在归一化后约为$(0.994,0.110)$。

同桶不是经核实的相关关系，多轮平均也不是全量softmax注意力的精确重写。它近似的是由碰撞概率诱导的加权方式。行为哈希与在线预测解耦节省了重复工作，但桶的版本、历史窗口及嵌入变化仍需维护；大量事件落入同一桶时，细粒度差别可能被求和掩盖。

% Outline: 2.2.1.15
\subsubsection{PinnerFormer：多时点目标的长期兴趣表示}
\label{sec:rs-pinnerformer}

一份批量更新、供多次请求使用的用户向量，不能只按“紧接着会点什么”判断是否合适。\term{PinnerFormer}把训练监督扩展到未来一段时间中的多个正反馈，让同一表示概括窗口内的兴趣。它延长的是预测范围，尚未估计某次推荐行动对未来的影响\scite{rs-r78}

输入包括物品内容向量、动作类型、入口、时间与持续时长等特征。因果Transformer编码各个已知前缀，通过投影及$\ell_2$归一化得到用户向量；物品侧也将内容向量映射到同一空间。训练在多个历史位置$k$选择未来窗口$P_k$中的正物品$i^+$，用采样softmax将$\bm{h}_k$与这些物品对齐。忽略采样校正时，一对正样本的目标为
\begin{equation}
  \ell_{k,i^+}=-\log
  \frac{\exp(\langle\bm{h}_k,\bm{v}_{i^+}\rangle/\tau)}
  {\exp(\langle\bm{h}_k,\bm{v}_{i^+}\rangle/\tau)
    +\sum_{j\in J_k}\exp(\langle\bm{h}_k,\bm{v}_j\rangle/\tau)}.
  \label{eq:rs-pinnerformer-loss}
\end{equation}
温度$\tau>0$控制分数差的尺度。原方案结合批内与随机负例，处理正物品误作负例，并对采样概率作logit校正；多个时点和窗口正例共同提供监督，而非让一个末端状态独自承担所有未来目标。

假设一次户外浏览之后，未来七天内用户收藏水壶、帐篷和摄影背包。下一动作目标只监督紧邻事件，窗口目标则可让三个收藏都约束同一前缀的表示。三者在训练时是未来标签，生成该前缀表示时都不可见。服务时批量编码用户，将向量用于近邻候选或下游评分特征。

窗口越长，标签越可能包含多个任务或变化的需求；是否合适取决于向量的使用周期和目标行为。向量一天更新一次也可能错过刚发生的兴趣变化，因而可以与第~\ref{sec:rs-transact}~节的即时分支组合。预测未来窗口的参与行为与控制未来参与行为之间，仍隔着行动选择和因果证据。

% Outline: 2.2.1.16
\subsubsection{TWIN-V2：层次压缩的长历史选择}
\label{sec:rs-twin-v2}

一致度量仍需访问大量原始事件。\term{TWIN-V2}先把相似历史压成簇，再在簇层面检索和精细聚合，以控制生命周期历史的规模。原任务是短视频CTR：先依据播放完成比例划分行为，再递归聚类，并限制每个簇的最大大小，避免大量相似记录挤在一个过大的簇中\scite{rs-r79}

每个簇形成一个虚拟物品。数值属性通过簇内聚合概括，类别属性按组成保留为聚合特征，模型继续使用TWIN式的物品自身特征与交叉特征分解。在线两阶段共享簇级相关性，并根据簇大小调整注意；最后把选中的簇表示聚合为兴趣，接入CTR预测，以曝光点击训练。

压缩必须保留数量的作用。若簇$r$的相关性为$a_r$、大小为$n_r$，两阶段使用调整分数$a'_r=a_r+\log n_r$，精细权重因而正比于$n_r e^{a_r}$。两个相同相关性的簇分别含10个和1个视频时，权重为$10/11$和$1/11$，而不是各占一半。但数量不能恢复具体发生顺序。即使两个簇的平均时长相同，前者可能由一半完整播放、一半快速跳过组成，后者可能全部为中等时长，因而原方案先按播放完成情况分组。

在商品教学类比中，十次相近背包浏览可由一个簇代表，在线访问从十条变为一个簇；这一映射仅解释压缩机制，不把视频的播放完成字段硬套成商品点击字段。聚类和物品表示更新会改变簇，簇版本须与在线特征协调。扩大覆盖历史与保留细节形成取舍，不能从“纳入更多记录”推断压缩无损。

% Outline: 2.2.1.17
\subsubsection{LONGER：逐层压缩的长历史建模}
\label{sec:rs-longer}

检索和离线聚类都在进入主网络前缩小历史。\term{LONGER}把一部分压缩放进网络，使较长历史先参与计算，再逐层降低后续的序列规模。候选及全局特征进入全局词元，邻近行为通过TokenMerge合并，局部InnerTrans处理组内交互；首层让较少的查询读取较长历史，后续在缩短的序列上继续编码\scite{rs-r80}

可以用张量形状理解这种分工。原历史长$n$，相邻$g$个事件组成一组后，组表示的数量约为$n/g$。组内变换处理局部事件关系，跨组层则处理压缩后的表示；首层查询数为$m<n$时，查询—历史关系数量由$n^2$降为$mn$。这个计数仅说明注意关系的规模，合并后的宽度、投影和前馈成本仍需另计。

例如，将“登山包、水壶、相机、镜头”分成两个相邻组，局部计算仍可使用每组内两个事件，后续看到的是两个组级表示；若先按候选检索只保留相机和镜头，登山组则根本不会进入精细网络。两种压缩保留的信息不同，不能仅比较最终都是两个输入位置就认为等价。

模型通过候选响应目标训练合并、局部处理与后续序列编码，最终产生用于CTR判断的表示和分数。因果注意力限制各位置读取历史的范围，全局词元如何放置也必须满足相应可见性约定。这里的“因果”指计算掩码，和通过干预识别推荐效果的因果推断含义不同。

合并后的组表示一般不能反解每条原始记录，局部顺序也可能只以压缩形式保留。缓存能否复用取决于候选是否已经影响缓存位置、历史是否追加以及参数是否变更。因而长上下文能力要同时检查信息粒度、训练范围与服务预算，不能仅用最大输入长度衡量。

% Outline: 2.2.1.18
\subsubsection{TiSASRec：时间间隔条件下的序列建模}
\label{sec:rs-tisasrec}

“登山包、水壶、相机”发生在半小时内还是三个月内，可能对应不同的判断。位置编号只给出次序，不能区分这种节奏。\term{TiSASRec}在SASRec式历史编码中增加两两时间间隔，并在注意的键和值两侧分别加入位置与间隔表示\scite{rs-r101}

设请求前可见的最小非零时间间隔为$\delta_u$，截断上界为$r_{\max}$，间隔索引为
\begin{equation}
  r_{k,l}=\min\left(r_{\max},
     \left\lfloor\frac{|\tau_k-\tau_l|}{\delta_u}\right\rfloor\right).
  \label{eq:rs-tisasrec-time}
\end{equation}
全部时间戳相同时需为尺度规定兜底值；训练回放不能从请求后的事件估计$\delta_u$。查询位置$k$对历史位置$l$的注意logit使用物品键、位置键及$r_{k,l}$的间隔键，聚合值也增加独立的位置值与间隔值。因果和填充掩码继续限制可见位置。

假设时间戳为0、2、8小时，最小非零间隔为2小时，相对首项的缩放时间为0、1、4。若$r_{\max}=3$，两两查表索引为
\[
  \bm{R}=
  \begin{pmatrix}0&1&3\\1&0&3\\3&3&0\end{pmatrix}.
\]
末端对第二项读取间隔3，对自身读取0。若把末项改成4小时，最后一行改为$(2,1,0)$；物品顺序没变，但键和值中的时间信息都变了。

模型用各前缀的下一物品和采样负例作二元交叉熵训练，推断以末端表示和物品向量内积评分。为说明改变间隔怎样传到输出，假设物品值为$(1,0),(0.8,0.2),(0,1)$，两种时间输入产生权重$(0.2,0.2,0.6)$与$(0.2,0.5,0.3)$，暂把值侧附加项设零，则表示从$(0.36,0.64)$变为$(0.60,0.40)$，两个坐标方向候选的次序也随之反转。这只是给定参数结果的教学计算，时间变化并不必然造成该方向的变化。

按个人最小间隔缩放保留了相对节奏，也可能让绝对时长相差很大的两段历史拥有相同索引。截断则使大间隔共享表示。两项处理都需要按数据条件检验，时间权重本身不能证明过去行为对现在选择具有因果作用。

% Outline: 2.2.1.19
\subsubsection{STAN：时间与空间条件下的地点匹配}
\label{sec:rs-stan}

下一地点受到空间关系影响，只有顺序仍不够。\term{STAN}以用户、地点和签到时间为输入，分别建立历史事件之间、候选地点与历史之间的时空间隔。第一层注意更新轨迹中的各事件表示，第二层让候选与全部更新后的事件匹配，避免过早将整个轨迹压成一个向量\scite{rs-r102}

历史地点的距离按经纬度计算，时间使用事件差值；间隔通过连续编码或相邻离散表示的插值进入注意关系。候选—历史时间差在原设定中用目标访问时刻计算。因此，应用若由用户指定“今晚去哪里”，可以使用这个已知时刻；预测一个时间尚未知晓的下一签到时，不能直接读取数据集中未来签到的真实时间。后者需要重新定义可用输入，并重新训练和评价。

用一个明确指定目标时刻的教学例子检查矩阵。为便于手算，以局部平面千米坐标代替经纬度球面距离：三个历史地点为$(0,0),(3,0),(3,4)$，签到时刻为8、9、11时，候选地点为A$(0,4)$、B$(6,4)$，请求目标时刻为12时。两组历史关系为
\[
  \bm{D}_{\mathrm{history}}=
  \begin{pmatrix}0&3&5\\3&0&4\\5&4&0\end{pmatrix},\qquad
  \bm{T}_{\mathrm{history}}=
  \begin{pmatrix}0&1&3\\1&0&2\\3&2&0\end{pmatrix}.
\]
候选—历史距离两行分别为$(4,5,3)$与$(\sqrt{52},5,3)$，时间两行均为$(4,3,1)$小时。即使A、B距最后地点都为3千米，它们与更早轨迹的关系仍有差别。把第二次签到从9时改到10时，会改变历史时间矩阵，并使候选时间行变为$(4,2,1)$。

轨迹聚合把这些关系加入事件注意，候选匹配再把内容兼容性与候选时空关系结合，经注意及可学习聚合得到地点得分。为了只隔离距离输入的作用，假设一个简化评分头固定使用$s_i=-\sum_k d_{i,k}$，则A、B分别为$-12$和约$-15.21$；这是完整网络中的一种教学比较函数，不是STAN保证学习到的距离规则。实际模型以真实下一地点作正例，每步重新采样一定数量的非目标地点，用sigmoid交叉熵训练。服务时去掉训练采样器，对适用地点目录计算分数。

签到并不覆盖全部实际活动，GPS还可能有噪声；距离接近也不代表道路可达或正在营业。模型输出的是访问倾向，候选资格需另查。时间与空间条件改变输入和评价任务，不能把得分提高解释为自动满足真实出行要求。

% Outline: 2.2.1.20
\subsubsection{AGREE：群组请求下的成员偏好聚合}
\label{sec:rs-agree}

共同选一部电影时，个人推荐列表的交集可能为空，固定平均又未必能解释群组的实际选择。\term{AGREE}让成员权重随候选改变，结合成员聚合与群组自身嵌入预测共同交互，同时用个人交互训练共享表示\scite{rs-r103}

设群组$g$的成员集合为$U_g$，成员向量为$\bm{p}_u$、群组ID向量为$\bm{q}_g$。候选相关注意网络给出
\begin{equation}
  \alpha_{u,i}=
   \frac{\exp a_\theta(\bm{p}_u,\bm{v}_i)}
        {\sum_{v\in U_g}\exp a_\theta(\bm{p}_v,\bm{v}_i)},\qquad
  \bm{h}_g(i)=\sum_{u\in U_g}\alpha_{u,i}\bm{p}_u+\bm{q}_g.
  \label{eq:rs-agree}
\end{equation}
预测网络读取$\bm{h}_g(i)$、$\bm{v}_i$及二者逐元素乘积，产生群组得分；个人路径用$\bm{p}_u$替换群组表示。两条路径共享成员和物品表示，以个人和群组的正负交互共同训练。原工作使用回归式成对目标，要求正负得分差接近1，可写成$(s_{g,i^+}-s_{g,j}-1)^2$，个人样本同理。

假设两名成员向量为$(1,0)$、$(0,1)$，某候选的权重为0.75、0.25，成员聚合为$(0.75,0.25)$；再加群组嵌入$(0.1,0.2)$，得到$(0.85,0.45)$。群组ID项提供了成员聚合之外的参数贡献。候选改变后权重也可能改变，服务时必须重新计算，不能把这一结果直接当作独立群组塔用于全部候选。

注意权重由预测损失学习，不是成员真实的表决权。未见过的新群组也没有已训练的$\bm{q}_g$，需要明确初始化或去掉该项后的适配方案。共同选择是否照顾每个成员，则属于第~\ref{chap:rs-constraints}~章的目标与约束，不能从群组预测准确率直接得到。

% Outline: 2.2.1.21
\subsubsection{GroupIM：稀疏群组反馈下的偏好学习}
\label{sec:rs-groupim}

临时群组很少重复出现，为每个群组单独学习一个ID向量可能没有足够数据。\term{GroupIM}从成员各自的交互向量编码个人表示，再通过均值、最大池化或注意聚合构造群组表示。它利用成员关系辅助学习，而不要求每个新群组先拥有一份成熟ID嵌入\scite{rs-r104}

令$\bm{x}_u\in\{0,1\}^{|I|}$为成员交互向量，$\bm{p}_u=f_{\mathrm{enc}}(\bm{x}_u)$，$\bm{h}_g=f_{\mathrm{agg}}(\{\bm{p}_u:u\in U_g\})$。目录预测头输出$p_\theta(i\mid g)=\operatorname{softmax}(\bm{W}\bm{h}_g)_i$，以共同交互的多项似然训练。成员历史不能一律等权当作共同偏好，所以模型另设双线性判别器
\begin{equation}
  D(u,g)=\sigma(\bm{p}_u^\top\bm{W}_{D}\bm{h}_g).
  \label{eq:rs-groupim-discriminator}
\end{equation}
真实成员与采样非成员构成对比样本，二元交叉熵推动用户—群组表示保留成员信息；这是互信息估计的训练约束，不是直接观测到群组协商过程。

判别分数还有第二个用途：以$D(u,g)$加权成员历史对群组目录预测的损失贡献。把常数归一化因子省略，该项为
\[
  \ell_{\mathrm{member},g}
    =-\sum_{u\in U_g}D(u,g)\sum_{i\in I}x_{u,i}\log p_\theta(i\mid g).
\]
训练交替进行：固定判别器，用共同反馈及加权成员反馈更新推荐路径；再训练对比目标，同时更新判别器、编码器和聚合器。个人交互可预训练编码器，服务阶段则直接编码成员历史、聚合并计算目录得分，不必运行负成员判别。

假设教学编码器给两名固定成员向量$(1,0)$与$(0,1)$。第三人是$(1,0)$时，均值群组向量为$(2/3,1/3)$；第三人换成$(0,1)$时，变为$(1/3,2/3)$。取$\bm{W}_D$为单位矩阵，第一名成员的判别值随之从$\sigma(2/3)\approx0.661$变为$\sigma(1/3)\approx0.583$，其历史监督贡献也改变。若目录头为单位映射，两个物品方向的概率从约$(0.583,0.417)$交换为$(0.417,0.583)$。

该例连起了成员变化、表示、判别权重和目录分布，但不能证明某人实际说服了其他成员。成员没有有效历史，或者群组内部存在强烈冲突时，辅助约束也无法凭空确定共同选择。表~\ref{tab:rs-behavior-extensions}把这些补充路线放回输入和输出接口，便于与前面的模型比较。

\begin{table*}[htbp]
  \centering
  \small
  \renewcommand{\arraystretch}{1.15}
  \begin{tabularx}{\linewidth}{@{}lXXXX@{}}
    \toprule
    模型 & 额外信息或改动 & 候选参与位置 & 中间输出 & 监督与应用输出 \\
    \midrule
    ETA & 长历史二进制检索 & 汉明距离及精细注意 & 选中事件的兴趣 & 曝光点击；CTR \\
    SDIM & 多轮哈希桶聚合 & 桶选择 & 跨轮聚合兴趣 & 曝光点击；CTR \\
    PinnerFormer & 多时点未来窗口标签 & 用户编码后 & 批量用户向量 & 窗口内正反馈；目录得分 \\
    TWIN-V2 & 播放信息与层次簇 & 簇检索及聚合 & 簇级长期兴趣 & 曝光点击；CTR \\
    LONGER & 局部合并、逐层压缩 & 全局词元 & 压缩序列表示 & 曝光点击；CTR \\
    TiSASRec & 真实时间间隔 & 历史编码后 & 末端历史向量 & 后继正负例；目录分数 \\
    STAN & 地点与时空间隔 & 轨迹匹配 & 各事件的更新表示 & 下一签到；地点得分 \\
    AGREE & 成员及群组ID & 成员聚合 & 候选相关群组向量 & 个人与共同交互；候选分数 \\
    GroupIM & 成员历史与关系 & 群组编码后 & 群组向量 & 共同／成员交互及对比；目录分布 \\
    \bottomrule
  \end{tabularx}
  \caption{长历史、情境和群组条件下的补充模型。相同的历史长度不意味着相同信息；同样输出目录得分，也不意味着具有相同的候选访问成本。}
  \label{tab:rs-behavior-extensions}
\end{table*}

% Outline: 2.2.2
\subsection{对象信息理解}
\label{sec:rs-item-understanding}

用户表示再丰富，也无法替代对象事实。推荐一个通勤背包至少需要知道价格、容量、重量和防水说明；只有标题里的“户外”一词，既不能确认容量，也不能确认防水等级。对象理解把文本、图片、视频和元数据变成可使用的信息，同时保留哪些是直接登记、哪些由模型推断、哪些尚缺依据。

不同模态提供的证据范围不同。图片可以展示颜色、外观和可见结构，却通常不足以确认准确重量；结构化参数可能明确重量，却需要检查版本、计量单位及是否包含配件。语言和视觉编码器可以提取特征，预测价值仍须由当前任务检验。某个向量能识别背包主题，不表示它保留了区分通勤和重装登山所需的全部细节。

对象之间的关系也要标明含义。两个背包外形相似，不一定互相替代；相机与适配存储卡可能互补，但买相机和买水壶经常共现，也可能只是同一批旅游用户的消费习惯。品牌和主题关系来自登记知识，共现关系来自交互，视觉相似来自内容计算。关系能否承担兼容性、替代性或组合要求，取决于证据是否支持这种用途。

因此，对象资料不应只有一个无来源的“特征值”。对于重量0.8千克，可以保存厂家规格和更新时间；对于“适合通勤”，可以保存判定依据及不确定性；遇到详情页和结构化字段冲突，应保留冲突状态而不是由模型分数覆盖。下节将内容与关系真正接到匹配目标；第~\ref{chap:rs-presentation}~章还会检查展示表达有没有超出这些事实。

% Outline: 2.2.3
\subsection{需求与对象匹配}
\label{sec:rs-user-item-match}

匹配包含两类问题：对象是否满足条件，以及在满足条件的对象中是否符合个体偏好。价格不超过500元是可以查字段的条件，偏爱简洁外观是需要结合个人信息的倾向。二者可以联合计算，但不能互相抵消：再高的外观偏好分数也不能证明600元的商品满足500元上限。

未知尤其不能被写成满足。假设三个背包价格为450、480、550元，重量分别为0.8千克、未知、0.7千克；用户要求价格不超过500元、重量不超过1千克。第一个通过已知条件，第二个待核验，第三个不通过价格条件。若语义匹配分数依次为0.7、0.9、0.95，分数也没有改变上述状态。搜索会将这类条件进一步联系到查询相关性和资格核验，见第~\ref{chap:rs-search}~章。

偏好匹配则需要把异质证据放到可比较的表示中。标题相近、交互相近和个体喜欢是不同关系，模型通过训练目标学习它们在当前数据中的对应。下面比较协同与文本对齐、知识关系传播、多模态构图和语言内容编码；这些是可替代或组合的建模选择，不是每次请求必须依次经过的步骤。

% Outline: 2.2.3.1
\subsubsection{A-LLMRec：协同与文本表示对齐}
\label{sec:rs-a-llmrec}

协同推荐器知道哪些物品在交互中相关，语言模型却未必能直接读取它的嵌入。\term{A-LLMRec}通过两阶段映射连接两种表示：先对齐协同与文本，再把联合表示投到语言模型的输入空间，用标题生成完成给定协议下的推荐\scite{rs-r23}

第一阶段固定预训练协同推荐器。物品协同向量$\bm{v}_i$和文本编码$\bm{t}_i$分别经过编码器$f_I,f_T$，用均方差使$f_I(\bm{v}_i)$接近$f_T(\bm{t}_i)$；各自的解码器重建原向量，避免只靠把所有表示压成同一点降低匹配误差。再用重建后的物品向量与协同用户向量进行正负物品预测。总目标为
\begin{equation}
  \ell_1=\ell_{\mathrm{match}}
      +\alpha\ell_{\mathrm{item-recon}}
      +\beta\ell_{\mathrm{text-recon}}+\ell_{\mathrm{rec}}.
  \label{eq:rs-a-llmrec-stage1}
\end{equation}
原方法在这一阶段微调作为文本编码器的Sentence-BERT；“协同推荐器冻结”不能被扩写成全部语言组件都冻结。

第二阶段固定前阶段使用的联合编码器和生成语言模型，训练用户、物品两个投影网络，将表示转换成语言模型词元维度。提示包含用户表示、历史标题及其表示、候选标题及其表示；目标是下一物品标题的逐词元负对数似然。因而更新的是连接网络，不是生成模型全部参数。

已有交互的登山包可走协同编码$f_I(\bm{v}_i)$；全新登山包没有训练过的协同ID向量时，可由文本路径$f_T(\bm{t}_i)$取得相同维度的联合表示。假设两条路径分别输出$(0.8,0.2)$和$(0.75,0.25)$，对齐目标缩小了表示差异，但不保证新包满足重量限制，也不证明两件物品对每个人同样合适。

模型生成标题后，仍需将输出映射到真实目录。重名、别名和目录外文本都要处理；原提示中的候选列表及实验协议也限制了它实际比较哪些物品。标题生成表现不能直接作为完整目录召回能力的证据。

% Outline: 2.2.3.2
\subsubsection{KGAT：知识与协同关系的注意力匹配}
\label{sec:rs-kgat}

用户只浏览过一个背包，但这个背包的品牌和主题还能连接其他对象。\term{KGAT}把用户—物品交互与物品知识三元组组织成统一图，既学习实体在关系中的匹配，也沿关系传播邻居信息，最终计算用户—物品偏好分数\scite{rs-r84}

知识三元组$(h,r,t)$表示头实体经关系$r$连接尾实体。关系投影矩阵$\bm{W}_r$与关系向量$\bm{e}_r$给出能量
\begin{equation}
  g(h,r,t)=\lVert\bm{W}_r\bm{e}_h+\bm{e}_r
                         -\bm{W}_r\bm{e}_t\rVert_2^2.
  \label{eq:rs-kgat-energy}
\end{equation}
能量越小，三元组在该表示中越相容。以真实三元组和替换尾实体得到的对照三元组作成对训练。邻居传播另用关系相关logit
$a(h,r,t)=(\bm{W}_r\bm{e}_t)^\top
\tanh(\bm{W}_r\bm{e}_h+\bm{e}_r)$，在头实体邻域归一化后加权汇总尾实体。

新节点表示结合自身与邻域。原模型比较相加、拼接和双交互聚合，后者同时变换$\bm{e}_h+\bm{e}_{N(h)}$与$\bm{e}_h\odot\bm{e}_{N(h)}$。多层传播取得多跳关系，拼接各层用户和物品表示后计算内积；训练交替使用知识三元组目标和BPR推荐目标。

在“用户—背包—品牌—另一背包”的小图上，多层传播可让用户表示收到品牌及关联对象的信息。若某节点的两个邻居权重为0.75、0.25，向量为$(1,0)$、$(0,1)$，邻域汇总便为$(0.75,0.25)$，再进入节点聚合，而不是直接成为推荐概率。关系必须已经取得并完成实体链接；错误品牌边也会进入传播。注意权重可以追踪计算依赖，不能被视为某品牌导致购买的因果证明。

% Outline: 2.2.3.3
\subsubsection{LATTICE：多模态物品图的偏好匹配}
\label{sec:rs-lattice}

有些物品关系没有人工登记，却能由图像或文本发现。\term{LATTICE}按各模态特征建立稀疏物品图，学习特征投影以调整邻域，再传播物品ID表示，将结果与协同表示结合用于推荐。内容在这里主要决定边，而不是直接取代全部协同计算\scite{rs-r85}

对于模态$m$，由原特征的余弦相似度构建初始近邻图$\bm{A}^{(m)}_0$；可学习投影产生新的相似图$\widetilde{\bm{A}}^{(m)}$。经过稀疏化和归一化后，混合初始关系与学习关系，再以可学习模态权重融合：
\begin{equation}
  \bm{A}=\sum_m\pi_m
   \bigl[\lambda\bm{A}^{(m)}_0
          +(1-\lambda)\widetilde{\bm{A}}^{(m)}\bigr],
  \qquad \sum_m\pi_m=1.
  \label{eq:rs-lattice-graph}
\end{equation}
沿融合图传播的是物品ID嵌入。将图分支表示加入协同模型的物品表示，与用户表示形成得分，以BPR目标联合训练模态投影、融合及推荐参数。

假设两款外观相近的背包在视觉图中的边权为0.9，文本用途差异使文本边权仅为0.2。暂不考虑图归一化变化，若视觉、文本权重为0.25、0.75，融合权重为0.375；改成0.75、0.25，则为0.725。这个局部算例说明模态取舍怎样改变传播，并没有证明较大的权重更正确。学习仍需看最终偏好监督是否支持该邻域。

相似图可能把旅行箱外观和通勤背包联系起来，却没有核验用途或互补性。新物品也需要可用内容、图节点及表示更新路径；在已有目录图上评测无交互节点，与无限接收训练后新节点是不同条件。模型输出匹配分数，不能代替属性事实检查。

% Outline: 2.2.3.4
\subsubsection{Recformer：物品文本与历史的联合匹配}
\label{sec:rs-recformer}

ID嵌入把物品绑定到某个目录；迁移到新目录时，即使两件商品描述相似，也没有自然共享的ID。\term{Recformer}直接把键值属性文本作为物品表达，再把历史物品文本组成用户输入，以共享语言编码器学习可比较的历史与物品向量\scite{rs-r86}

输入不只是把若干标题无分隔地连接起来。词元嵌入之外，类型、词元位置和物品位置嵌入区分属性结构及行为归属；Longformer在局部注意窗口及全局位置间交换信息，得到历史表示$\bm{h}_u$。同一编码结构处理单个物品文本，得到$\bm{v}_i$，用温度缩放的余弦相似度预测下一物品。

预训练结合掩码语言任务和物品对比任务：前者恢复被隐藏的属性词元，后者将历史表示与真实下一物品拉近、与采样物品区分。推荐微调还要维护目录表示，原方案先随训练阶段更新物品表示表，再固定选定的物品表继续适配历史编码；若编码器已经改变而物品缓存完全未更新，匹配空间可能错位。

例如，历史背包文本包含“品牌：A；材质：尼龙；用途：通勤”，新物品包含“品牌：B；材质：尼龙；用途：通勤”。文本路径让两者能共享材质和用途信号。假设归一化历史向量为$(1,0)$，新包和相机向量分别为$(0.8,0.6)$与$(0,1)$，余弦得分为0.8和0。真实模型从文本和反馈学习这些向量，不是预先把通勤指定成第一个坐标。

服务时先编码当前目录，历史编码后通过相似度访问物品；这与自由生成标题不同。双向编码仅作用于已经发生的历史，不能读入未来目标的文本作为历史事件。文本遗漏关键属性或宣传词压过实际用途时，相似性仍可能失准，需要单独检验内容质量与条件符合情况。

% Outline: 2.2.3.5
\subsubsection{HLLM：语言内容与行为的分层匹配}
\label{sec:rs-hllm}

把所有历史物品的完整文本交给同一个编码器，会反复处理多人共同浏览的物品描述。\term{HLLM}将物品内容编码和用户行为编码分为两层：Item LLM从物品文本及专用词元取得向量，User LLM把这些向量作为序列输入，学习下一物品或候选响应\scite{rs-r87}

对物品文本$T_i$，第一层给出$\bm{v}_i=f_{\mathrm{item}}(T_i)$；第二层读取请求前的$\bm{v}_{i_1},\ldots,\bm{v}_{i_n}$，形成用户输出$\bm{h}_{u,t}$。下一物品任务用实际后继和随机负例构成InfoNCE目标，即提高正物品在正负集合内的归一化相似度，其形式与式~\eqref{eq:rs-pinnerformer-loss}相近，但监督是后继物品，不是同一个未来窗口任务。训练更新参与适配的两层模型。

候选响应任务有两种接口。早融合把候选向量追加到用户序列，让候选参与深层计算，再接点击判别头；后融合先独立编码历史和候选，在输出侧进行交互。后融合可以复用同一用户表示，早融合需要保留各候选的相关计算。二者以点击交叉熵训练，还可结合下一物品辅助目标；不能因共用Item LLM就忽略使用接口差异。

假设100名用户都浏览同一款背包，冻结物品编码器后，其内容向量可以先算一次，再由100段不同历史复用。若下一物品输出为$(0.6,0.8)$，背包与相机的归一化物品向量分别为$(1,0)$、$(0,1)$，相似度为0.6和0.8；这只是候选匹配结果，不是模型生成了一段相机介绍。物品文本改变或编码器继续更新时，缓存也须随版本刷新。

层次编码减少重复内容处理，不保证内容向量保留一切细节，也不能把原HLLM的向量预测扩写成广告创意生成。文本事实、训练时可用历史和目录映射仍须各自核验。

% Outline: 2.3
\section{信息获取与维护}
\label{sec:rs-evidence-maintenance}

已有信息足够时，问题是怎样编码；缺少或过时的信息则需要先处理证据本身。补充提供新的依据，更新让依据反映当前状态，检验确定依据能支持什么结论。这三项职责贯穿用户、对象和匹配建模，不能等评分器训练完才统一补做。

% Outline: 2.3.1
\subsection{缺失信息补充}
\label{sec:rs-missing-information}

“\term{冷启动}”常把不同缺口混在一起。新用户可能没有历史却愿意填写用途；匿名用户可能有当前会话却没有稳定身份；新物品可能没有交互却有完整属性；长期低曝光物品可能存在很久，仍缺少可比较的反馈。训练时未见实体又是模型参数覆盖条件，和有没有内容不是同一回事。先指出缺什么，才能选择内容先验、显式初始化、关系共享或少量反馈适应。

已有其他业务的信息时，还需建立对应关系。跨域推荐涉及不同目录或交互分布，至少要说明用户或物品是否重叠、迁移单向还是双向、目标域零反馈还是少反馈。同一商品目录在首页和购物车的CTR适配是多场景问题，留到第~\ref{chap:rs-ranking}~章讨论。不能因为都叫“迁移”，就默认身份、对象和行为含义已经对齐。

能询问用户时，补充信息也可以直接来自回答。例如，不清楚“轻便”指重量上限还是肩带舒适度，询问后可以改变可选集合；若两个解释最终产生完全相同的选择，继续询问的收益可能很小。问题价值要和打断成本一起衡量，回答还需带上适用范围。用户此次答“买给同事”，不应被永久登记为本人偏好改变。交互界面与展示测试分别由第~\ref{chap:rs-presentation}~章承接。

% Outline: 2.3.1.1
\subsubsection{DropoutNet：偏好信息缺失下的内容补充}
\label{sec:rs-dropoutnet}

已有潜在模型在热物品上有效，但新物品没有学好的交互向量。\term{DropoutNet}从这个模型出发，把用户和物品的偏好向量与内容特征分别送入映射网络；训练时随机移除整侧偏好输入，让网络学习在该输入缺失时依靠内容恢复原模型的相关性得分\scite{rs-r02}

设原向量为$\bm{p}_u,\bm{q}_i$，内容为$\bm{c}_u,\bm{c}_i$，随机掩码$m_u,m_i\in\{0,1\}$控制偏好侧输入。新表示及目标可写成
\begin{align}
  \widehat{\bm{p}}_u&=f_U([m_u\bm{p}_u;\bm{c}_u]),&
  \widehat{\bm{q}}_i&=f_I([m_i\bm{q}_i;\bm{c}_i]),\\
  \ell_{u,i}&=\bigl(
    \langle\widehat{\bm{p}}_u,\widehat{\bm{q}}_i\rangle
      -\langle\bm{p}_u,\bm{q}_i\rangle\bigr)^2.
  \label{eq:rs-dropoutnet}
\end{align}
有偏好输入的样本促使网络保留原模型的信息，偏好被移除的样本迫使内容参与恢复。丢弃的是整个偏好向量，不等于任意隐藏层的常规dropout。

假设原模型给一位用户和旧登山包的得分为0.8，训练将包的偏好向量置零，内容网络给出的新内积为0.6，则平方误差为0.04，梯度推动内容路径解释缺失的部分。新登山包服务时同样以零偏好向量和真实属性输入。新用户取得少量交互后，原方案还以交互过物品的表示构造偏好输入近似，并在训练中模拟这种早期状态。

输出是新向量与匹配分数，能够接入内积召回；其监督目标仍来自原潜在模型。原模型有偏差、内容无信息或内容与交互的对应关系改变时，恢复教师得分也不保证恢复真实偏好。两侧都缺少可用信息时，输入缺失模拟不能凭空补出依据。

% Outline: 2.3.1.2
\subsubsection{KAR：语言知识与行为的适应匹配}
\label{sec:rs-kar}

商品资料可能未覆盖领域中的全部选择因素，行为记录也未直接说明用户为何选择。\term{KAR}按领域因素组织提示，让语言模型离线生成用户偏好推理文本和物品知识文本，再将这些文本编码并适配为推荐特征\scite{rs-r88}

偏好提示以已有历史为条件，知识提示以物品信息及领域因素为条件。离线文本编码得到$\bm{t}_u,\bm{t}_i$，混合专家适配器包含共享专家和分别服务用户推理、物品知识的专属专家。用户侧可概括为
\[
  \bm{z}_u=\sum_{e\in E_{\mathrm{shared}}\cup E_{\mathrm{user}}}
                \alpha_e(\bm{t}_u)f_e(\bm{t}_u),
\]
物品侧同理但使用自己的专属专家。门控权重选择各专家贡献；适配输出与既有用户、物品及情境字段一起进入推荐网络，随点击等下游目标训练。语言生成和编码可预先完成，服务时读取特征。

假设背包知识补充了容量、面料和使用场景三个方面，适配器中共享与专属专家的教学输出分别为$(0.6,0.2)$和$(0.2,0.8)$，门权重为0.25和0.75，则输出为$(0.3,0.65)$。这个向量只是评分的新输入，尚未作出推荐。若生成文本错误地把普通防泼水写成可完全浸水，预测网络也不会自动把它纠正成真实规格。

因此，物品事实需要核验来源，偏好推理应保留为待检验假设。丰富的描述和流畅的理由不构成用户真实动机的证据，更不构成因果解释。相较第~\ref{sec:rs-a-llmrec}~节的表示对齐，KAR主要增加外部生成内容，二者信息来源及验证责任不同。

% Outline: 2.3.1.3
\subsubsection{MeLU：少量反馈的用户偏好适应}
\label{sec:rs-melu}

有了几次真实反馈，仍可能不足以为新用户独立训练完整网络。\term{MeLU}从已有用户学习一份容易适应的初始化：把每个用户作为一项训练任务，用少量支持交互调整预测层，再用其他交互检查调整是否有用\scite{rs-r89}

把参数分为属性编码部分$\theta_{\mathrm{emb}}$和偏好预测部分$\theta_{\mathrm{pred}}$。对用户$u$的支持集$S_u$进行局部更新：
\begin{equation}
  \theta'_{\mathrm{pred},u}=
  \theta_{\mathrm{pred}}-\alpha\nabla_{\theta_{\mathrm{pred}}}
     \ell(S_u;\theta_{\mathrm{emb}},\theta_{\mathrm{pred}}).
  \label{eq:rs-melu-inner}
\end{equation}
此时保持属性嵌入部分不变。随后在该用户查询集$Q_u$上评价适应后的预测，跨用户最小化
$\sum_u\ell(Q_u;\theta_{\mathrm{emb}},\theta'_{\mathrm{pred},u})$，
更新共享初始化及属性编码。支持集与查询集承担不同训练职责，不能用同一反馈反复证明适应后泛化更好。

用标量预测层说明一步适应。假设背包特征$x=1$、评分$y=4$、初始参数$w=3$，损失为$\tfrac12(wx-y)^2$，步长为0.2。梯度为$-1$，更新后$w'=3.2$；特征$x=0.5$的候选预测从1.5变为1.6。真实MeLU以用户和物品属性进入多层预测器，这里仅隔离局部梯度怎样改变后续判断。

服务时根据新用户少量标注执行局部更新，再对候选评分。完全没有反馈时，式~\eqref{eq:rs-melu-inner}无从执行，只能使用共享先验。元训练用户与新用户差异很大、支持反馈噪声较多时，快速适应也可能快速过拟合；取得哪些反馈仍是独立的信息获取问题。

% Outline: 2.3.1.4
\subsubsection{MetaEmb：新广告嵌入的内容初始化与反馈适应}
\label{sec:rs-metaemb}

新广告已经有图片、类别和商品属性，却没有成熟的ID向量。\term{MetaEmb}学习一个由内容产生初始嵌入的生成器，并同时考察它在零反馈时的预测和少量反馈更新后的预测。与MeLU主要适应用户预测层不同，这里适应的是新广告的ID嵌入\scite{rs-a28}

固定已有CTR主网络参数$\theta$，生成器$g_\phi$由广告特征$\bm{c}_i$给出$\bm{v}^{(0)}_i=g_\phi(\bm{c}_i)$。在旧广告上模拟新ID，取两组少量样本$D_i^a,D_i^b$，计算
\begin{align}
  \bm{v}^{(1)}_i&=\bm{v}^{(0)}_i
      -\eta\nabla_{\bm{v}^{(0)}_i}\ell(D_i^a;\theta,\bm{v}^{(0)}_i),\\
  \ell_{\mathrm{meta},i}&=
     \lambda\ell(D_i^a;\theta,\bm{v}^{(0)}_i)
      +(1-\lambda)\ell(D_i^b;\theta,\bm{v}^{(1)}_i).
  \label{eq:rs-metaemb}
\end{align}
生成器通过两项损失学习：初始向量应立即可用，也应在少量样本更新后表现合理。主CTR网络在这一训练过程中冻结，不能把其变化误记为初始化效果。

假设某新背包广告的初始标量嵌入为0.2，教学CTR头为$\sigma(v)$。收到一次点击后，交叉熵对$v$的梯度为$\sigma(0.2)-1\approx-0.450$；步长0.1时，更新到约0.245，预测由0.550变为约0.561。真实投放中的其他字段和多条样本会改变这个结果，算例仅区分内容初始化与反馈适应两步。

没有反馈时只执行初始化，反馈到来后才更新ID向量。推广到普通新物品时需要明确CTR任务和字段的对应；广告资格、支付和预算都不由这个嵌入生成器决定。内容缺失及新广告分布变化也会限制初始化的可靠性。

% Outline: 2.3.1.5
\subsubsection{EMCDR：跨域表示映射的偏好补充}
\label{sec:rs-emcdr}

同一用户在来源目录有充足交互，在目标目录却没有记录，可以尝试把已知偏好映到目标空间。\term{EMCDR}分别在两域用矩阵分解或BPR学习潜在因子，再用重叠用户或物品训练映射。两域因子学习和跨域映射是两个阶段\scite{rs-r95}

以重叠用户集合$U_{\cap}$为例，来源和目标因子为$\bm{p}_u^S,\bm{p}_u^T$，映射$f_\phi$通过
\begin{equation}
  \min_\phi\sum_{u\in U_{\cap}}
     \lVert f_\phi(\bm{p}_u^S)-\bm{p}_u^T\rVert_2^2
  \label{eq:rs-emcdr-map}
\end{equation}
学习，$f_\phi$可以是线性映射或MLP。对仅在目标侧缺交互的用户，以$f_\phi(\bm{p}_u^S)$和目标域已有物品因子计算内积，再进入目标目录的Top-$K$访问。重叠物品可建立相应的物品映射任务，但不是所有应用都会同时具有两种重叠。

假设来源用户因子为$(1,0)$，线性映射为$\operatorname{diag}(2,1)$，目标用户表示为$(2,0)$。目标候选向量为$(1,1)$与$(0,2)$，得分分别为2和0。映射改变的是表示坐标及关联，不表示来源域的评分2与目标域评分2拥有相同原始含义。

这个过程需要来源因子、对齐实体和可学习的跨域关系。没有重叠实体，式~\eqref{eq:rs-emcdr-map}就缺少监督；两域偏好关系变化时，拟合良好的旧映射也可能造成负迁移。只能取得聚合数据或身份使用受限时，需要另建可用证据条件，不能默认原方法仍可直接应用。

% Outline: 2.3.1.6
\subsubsection{CoNet：跨域网络连接的偏好共享}
\label{sec:rs-conet}

最终因子的映射只在一个位置交换信息。\term{CoNet}让两个目录的预测网络在多层计算中共享隐藏信息。原设定以相同用户连接两域，各域保留独立物品表和预测头，通过跨域连接联合训练\scite{rs-r96}

对两个域$S,T$，一对连接层可写为
\begin{align}
  \bm{h}^{S,(\ell+1)}
    &=\rho(\bm{W}_S^{(\ell)}\bm{h}^{S,(\ell)}
           +\bm{H}_{T\to S}^{(\ell)}\bm{h}^{T,(\ell)}
           +\bm{b}_S^{(\ell)}),\\
  \bm{h}^{T,(\ell+1)}
    &=\rho(\bm{W}_T^{(\ell)}\bm{h}^{T,(\ell)}
           +\bm{H}_{S\to T}^{(\ell)}\bm{h}^{S,(\ell)}
           +\bm{b}_T^{(\ell)}).
  \label{eq:rs-conet-cross}
\end{align}
$\rho$为激活函数；$\bm{H}$传递另一域当前层的信息，$\bm{W}$处理本域信息。原CoNet令两个方向共享同一矩阵，即$\bm{H}_{T\to S}^{(\ell)}=\bm{H}_{S\to T}^{(\ell)}$，这里的方向下标只是区分两次调用。两域的真实交互及采样未观测物品提供各自的二元交叉熵，共同更新共享用户因子、域内网络与连接；SCoNet再对连接矩阵施加稀疏约束，以减少不必要的信息传递。

假设本域变换和偏置合计0.3，跨域贡献为0.2，则ReLU输出0.5；把连接置零，输出回到0.3。这可作为局部消融的教学计算，实际负迁移应比较目标域完整预测，而不能以非零连接数量判断。服务时还需按训练定义提供跨域分支的可用输入，再由目标域预测头比较物品；共享一个用户因子并没有消除候选相关网络。

模型需要原设定中的用户对应和两域训练信息。连接稀疏不保证迁移方向正确，也不覆盖任意无重叠目录或目标域零反馈任务。若只有来源历史而无法构造训练时使用的另一域输入，应重新设计接口并评价，不能把缺失分支静默当作已经共享的知识。

% Outline: 2.3.2
\subsection{理解依据更新}
\label{sec:rs-evidence-refresh}

信息缺失与信息过时需要不同处理。用户刚从登山改为通勤需求，系统可能已经收到明确查询却仍使用旧画像；也可能查询已进入输入，但模型仍过度依赖旧历史。前者是证据接入延迟，后者是计算或学习对变化适应不足，增加更多旧记录未必有帮助。

固定请求发生在10:05，假设用户10:04收藏相机，长期向量上次更新在午夜，模型参数更新在9:00。TransAct若在10:05读到该收藏，可以改变即时分支；长期向量仍未更新并不妨碍这一点，参数也不必在一分钟内重新训练。相反，即使9:00的模型比昨日模型更准确，它也不能使用尚未到达的10:04事件。

MIMN需要检查行为是否已写入记忆，SIM需要检查事件是否已进入可检索历史，TWIN还需检查物品投影与检索参数版本。四种状态不能统一成一个“更新时间”。图~\ref{fig:rs-history-access}比较事件保存和请求读取的差异；信息路径改变之后，延迟、存储及重算成本也会改变。

\begin{figure*}[htbp]
  \centering
  % 三条历史访问路径；最终宽度约162mm，按保存粒度比较。
\begin{tikzpicture}[
  font=\figurefont\footnotesize,
  rsbox/.style={rectangle,draw=BookRule,fill=BookPaper,
    text=BookInk,align=center,text width=40mm,minimum height=13mm,inner sep=4pt},
  rsflow/.style={-{Stealth[length=2mm]},draw=BookMuted,line width=.9pt}
]
  \node[font=\figurefont\footnotesize\bfseries] at (0,1.2) {新事件到来};
  \node[font=\figurefont\footnotesize\bfseries] at (5.7,1.2) {保留的信息};
  \node[font=\figurefont\footnotesize\bfseries] at (11.4,1.2) {预测请求读取};
  \node[rsbox,fill=FigureInput!10] (a1) at (0,0) {增量读写\\擦除、添加、状态更新};
  \node[rsbox,fill=FigureModel!10] (b1) at (5.7,0) {固定容量记忆\\部分事件细节被压缩};
  \node[rsbox,fill=FigureOutput!10] (c1) at (11.4,0) {读取记忆及归纳状态\\结合候选预测};
  \node[rsbox,fill=FigureInput!10] (a2) at (0,-1.9) {追加原始事件\\更新历史索引};
  \node[rsbox,fill=FigureModel!10] (b2) at (5.7,-1.9) {可检索事件记录\\保留逐条行为信息};
  \node[rsbox,fill=FigureOutput!10] (c2) at (11.4,-1.9) {候选检索相关事件\\再作精细聚合};
  \node[rsbox,fill=FigureInput!10] (a3) at (0,-3.8) {更新哈希桶或行为簇\\刷新聚合及版本};
  \node[rsbox,fill=FigureModel!10] (b3) at (5.7,-3.8) {桶或簇表示\\保留聚合信息与数量};
  \node[rsbox,fill=FigureOutput!10] (c3) at (11.4,-3.8) {候选访问桶或簇\\形成聚合兴趣};
  \foreach \r in {1,2,3} {
    \draw[rsflow] (a\r.east) -- (b\r.west);
    \draw[rsflow] (b\r.east) -- (c\r.west);
  }
\end{tikzpicture}

  \caption{长历史的三种访问路径。增量记忆在新事件到来时压缩状态；历史检索在请求时按候选选事件；桶或簇聚合预先改变历史粒度，再由请求取得相关聚合。图比较信息保存与访问职责，不表示各模型的全部预测网络相同。}
  \label{fig:rs-history-access}
\end{figure*}

更新评价应固定请求时点，记录各信息真正可用的时间，比较变化发生后的预测质量和维护代价。目录可用性更新由第~\ref{chap:rs-retrieval}~章承接，模型重训与多场景适应由第~\ref{chap:rs-ranking}~章展开，具体状态同步属于系统卷。更及时的信息可以帮助当前判断，但是否改善长期用户收益仍需后续决策和评价证据。

% Outline: 2.3.3
\subsection{理解质量检验}
\label{sec:rs-understanding-quality}

理解质量不能由画像字数或历史长度证明。用户侧应检查判断是否符合明确表达、近期变化及真实反馈；对象侧应检查属性和关系是否有来源、是否与当前版本一致；匹配侧应检查条件符合情况和偏好预测是否各自成立。生成画像尤其要能回到原始证据，不能让生成内容在后续链条中变成无来源的事实。

比较模型时，先固定证据再改变处理方式。例如，在相同历史上移除顺序信息，检验顺序的贡献；固定事件后改变会话边界，检验划分规则；固定完整历史后改变检索条数，检验选择损失。再单独改变历史长度或信息时效，观察额外证据的贡献。若一次同时增加历史、内容和计算量，最终改善无法归到某一种编码机制。

分组检查应覆盖新老用户、稀疏与密集交互、稳定与匿名身份、不同证据来源及变化后的阶段。没有已知事实时，正确输出“未知”可能比给出自信但错误的属性更有用；取得补充信息的成本也要计入。成员注意、图注意和兴趣状态可用于检查网络依赖了什么，但不能独立验证心理或因果解释。

本章检验的是当前依据是否足够、是否一致及是否及时。训练标签如何校正曝光、位置和选择偏差，由第~\ref{chap:rs-ranking}~章建立；完整数据切分、指标与效果识别由第~\ref{chap:rs-evaluation}~章讨论。局部模型比较应保留这些入口，不能以一次离线命中提升结束全部验证。

% Outline: 2.4
\section{本章小结}
\label{sec:rs-understanding-summary}

从登山历史推断通勤请求，关键是分清已知表达、可见行为、对象事实和模型推断。直接条件规定本次需求的一部分，历史和情境补充尚未表达的部分；对象表示提供属性与关系，匹配模型再学习这些证据如何共同支持选择。语义相似和高偏好分数都不能替代硬条件核验，未知信息也应保留为未知。

行为模型在顺序、会话、候选参与、记忆容量和监督时间范围上作出不同选择。它们形成的中间表示及最终输出决定了下一环节怎样使用：候选无关向量适合共享编码，候选相关聚合需要保留逐候选计算，群组请求还需成员关系和共同标签。内容补充、少样本适应及跨域共享各有不同的信息前提，更新与质量检验则限定这些依据在当前请求中是否有效。

现在已有了能够比较需求与对象的表示，但表示本身尚未从大目录中取得真实候选。下一章将把内积、邻域关系与生成输出变成具体访问方法，并检查目录覆盖、资格和候选遗漏。理解形成的证据只有进入这些实际选择，才能继续接受任务效果的检验。
